\chapter{Schwinger--Keldysh Hydrodynamic Effective Actions}
\label{ch:thermal-sk-hydrodynamic-effective-actions}

\ConventionsRealTime

Hydrodynamic Schwinger--Keldysh effective actions are real-time doubled
generating functionals for conserved densities in a long-wavelength state.
Their defining constraints are inherited from the closed time path:
normalization at coincident sources, a reality relation, positivity of the
imaginary part, gauge or diffeomorphism Ward identities, and the KMS condition
when the initial state is thermal.  This chapter develops the
quadratic charge-diffusion action explicitly, because it is the simplest
place where all of these constraints can be seen without hiding the
source-response signs.

Let spacetime be \(\mathbb R_t\times\mathbb R^d\).  Spatial indices
\(i,j\) are raised with \(\delta^{ij}\), and
\(\nabla^2=\partial_i\partial_i\).  Fourier modes are written as
\(\exp(-\ii\omega t+\ii\mathbf k\cdot\mathbf x)\) in the derivations below.
The microscopic retarded commutator convention is the one fixed in
Chapter~\ref{ch:thermal-spectral-kubo-transport}.  In this chapter
\(K^R\) denotes the source-response kernel obtained by differentiating the
Schwinger--Keldysh functional with respect to the \(a\)-type source.  The
conversion between \(K^R\) and the commutator retarded function includes the
sign convention of the microscopic source coupling and possible local contact
terms.

\section{Closed-Time-Path Constraints}

Let \(A_{1\mu}\) and \(A_{2\mu}\) be background gauge fields on the two
branches of the Schwinger--Keldysh contour.  For an initial density operator
\(\rho_0\), define
\begin{equation}
  Z[A_1,A_2]
  =
  \operatorname{Tr}\!\left(
  U[A_1]\rho_0U[A_2]^\dagger
  \right).
  \label{eq:sk-hydro-generating-functional}
\end{equation}
Introduce average and difference variables
\begin{equation}
  A_{r\mu}=\frac12(A_{1\mu}+A_{2\mu}),
  \qquad
  A_{a\mu}=A_{1\mu}-A_{2\mu}.
  \label{eq:sk-hydro-ra-sources}
\end{equation}
If \(Z=\exp(\ii I_{\rm SK})\), unitarity and Hermiticity imply
\begin{align}
  I_{\rm SK}[A_r,A_a=0]&=0,
  \label{eq:sk-hydro-unitarity-zero}\\
  I_{\rm SK}[A_r,A_a]^*
  &=
  -I_{\rm SK}[A_r,-A_a],
  \label{eq:sk-hydro-reality}\\
  \operatorname{Im} I_{\rm SK}[A_r,A_a]&\ge0
  \label{eq:sk-hydro-positivity}
\end{align}
for source configurations in the domain of the effective description.
Equation \eqref{eq:sk-hydro-unitarity-zero} says that every term in the
effective action contains at least one \(a\)-type field.  The positivity
condition is the damping condition in the path-integral weight
\(\exp(\ii I_{\rm SK})\).

\section{Hydrodynamic Phase Field and Gauge Invariance}

For a conserved \(U(1)\) charge, the hydrodynamic mode is represented by
phase fields
\[
  \varphi_1,\qquad \varphi_2 .
\]
The branchwise gauge transformations are
\[
  A_{s\mu}\mapsto A_{s\mu}+\partial_\mu\lambda_s,
  \qquad
  \varphi_s\mapsto\varphi_s-\lambda_s,
  \qquad s=1,2.
\]
Thus the gauge-invariant combinations are
\begin{equation}
  B_{s\mu}=A_{s\mu}+\partial_\mu\varphi_s .
  \label{eq:sk-hydro-branch-b-fields}
\end{equation}
In \(r/a\) variables,
\begin{equation}
  B_{r\mu}=A_{r\mu}+\partial_\mu\varphi_r,
  \qquad
  B_{a\mu}=A_{a\mu}+\partial_\mu\varphi_a .
  \label{eq:sk-hydro-ra-b-fields}
\end{equation}
The field \(\varphi_a\) is the Lagrange-multiplier partner of the
conservation law.  Varying the action with respect to \(\varphi_a\) is
equivalent to the Ward identity because \(\delta B_{a\mu}=\partial_\mu
\delta\varphi_a\).

\section{Quadratic Diffusion Action}

Consider a homogeneous thermal state at temperature \(T>0\), with charge
susceptibility \(\chi>0\) and conductivity \(\sigma\ge0\).  Define
\[
  D=\frac{\sigma}{\chi}.
\]
The leading quadratic local action for the diffusive mode is
\begin{equation}
  I_{\rm diff}
  =
  \int dt\,d^dx\,
  \left[
  \chi B_{a0}B_{r0}
  -\sigma B_{ai}\partial_t B_{ri}
  +\ii T\sigma B_{ai}B_{ai}
  \right],
  \label{eq:sk-hydro-diffusion-action}
\end{equation}
up to terms with more derivatives or more powers of \(a\)-fields.  The
current obtained by varying with respect to the \(a\)-type source at
\(A_a=\varphi_a=0\) is
\begin{equation}
  n=\frac{\delta I_{\rm diff}}{\delta A_{a0}}
  =\chi B_{r0},
  \qquad
  j^i=\frac{\delta I_{\rm diff}}{\delta A_{ai}}
  =-\sigma\partial_t B_{ri}.
  \label{eq:sk-hydro-density-current-from-action}
\end{equation}
Using \(B_{r0}=A_{r0}+\partial_t\varphi_r\) and
\(B_{ri}=A_{ri}+\partial_i\varphi_r\), this current may be written as
\begin{equation}
  j^i
  =
  -D\partial_i n+\sigma E_i,
  \qquad
  E_i=\partial_i A_{r0}-\partial_t A_{ri}.
  \label{eq:sk-hydro-constitutive-current-with-source}
\end{equation}
The sign of \(E_i\) is fixed by the source convention in which a static
positive \(A_{r0}\) increases the charge density by \(\chi A_{r0}\).

\begin{proposition}[Diffusion equation from the \(a\)-phase]
\label{prop:sk-hydro-diffusion-equation}
At \(A_a=\varphi_a=0\), varying
\eqref{eq:sk-hydro-diffusion-action} with respect to \(\varphi_a\) gives
\begin{equation}
  \partial_t n+\partial_i j^i=0,
  \label{eq:sk-hydro-continuity-from-phase}
\end{equation}
with \(n\) and \(j^i\) as in
\eqref{eq:sk-hydro-density-current-from-action}.  With external sources set
to zero this becomes
\begin{equation}
  \partial_t n-D\nabla^2 n=0 .
  \label{eq:sk-hydro-diffusion-equation}
\end{equation}
\end{proposition}

\begin{proof}
Under \(\varphi_a\mapsto\varphi_a+\delta\varphi_a\),
\[
  \delta B_{a0}=\partial_t\delta\varphi_a,
  \qquad
  \delta B_{ai}=\partial_i\delta\varphi_a .
\]
The terms linear in \(a\)-fields therefore vary as
\[
  \delta I_{\rm diff}
  =
  \int dt\,d^dx\,
  \left[
  \chi(\partial_t\delta\varphi_a)B_{r0}
  -\sigma(\partial_i\delta\varphi_a)\partial_tB_{ri}
  \right],
\]
because the quadratic \(B_aB_a\) term does not contribute at
\(A_a=\varphi_a=0\).  After integrating by parts and discarding boundary
terms compatible with the closed-time-path variational problem,
\[
  \delta I_{\rm diff}
  =
  -\int dt\,d^dx\,\delta\varphi_a
  \left[
  \partial_t(\chi B_{r0})
  -\partial_i(\sigma\partial_tB_{ri})
  \right].
\]
Using \eqref{eq:sk-hydro-density-current-from-action}, the Euler-Lagrange
equation is exactly \eqref{eq:sk-hydro-continuity-from-phase}.  If
\(A_{r\mu}=0\), then \(n=\chi\partial_t\varphi_r\) and
\[
  j^i=-\sigma\partial_i\partial_t\varphi_r
  =-D\partial_i n,
\]
which gives \eqref{eq:sk-hydro-diffusion-equation}.
\end{proof}

\section{Density Response Kernel}

The density response to a scalar source follows from the same saddle
equation.  Set \(A_{ri}=0\) and use a Fourier mode
\(\exp(-\ii\omega t+\ii\mathbf k\cdot\mathbf x)\).  The equation of motion
from Proposition~\ref{prop:sk-hydro-diffusion-equation} gives
\[
  \chi(A_{r0}-\ii\omega\varphi_r)+\sigma k^2\varphi_r=0,
\]
because \(n=\chi(A_{r0}-\ii\omega\varphi_r)\) and
\(\partial_i j^i=\sigma k^2(-\ii\omega)\varphi_r\).  Hence
\[
  \varphi_r
  =
  -\frac{\chi}{\sigma k^2-\ii\chi\omega}A_{r0}.
\]
Substituting into \(n=\chi(A_{r0}-\ii\omega\varphi_r)\) gives the
source-response kernel
\begin{equation}
  K^R_{nn}(\omega,\mathbf k)
  =
  \frac{\delta n}{\delta A_{r0}}
  =
  \chi\,\frac{Dk^2}{Dk^2-\ii\omega}.
  \label{eq:sk-hydro-density-response-kernel}
\end{equation}
It has the diffusive pole
\begin{equation}
  \omega=-\ii Dk^2 .
  \label{eq:sk-hydro-density-response-pole}
\end{equation}
At fixed nonzero \(k\), the static limit is
\[
  K^R_{nn}(0,\mathbf k)=\chi.
\]
At \(k=0\) and nonzero \(\omega\), the kernel vanishes, expressing
conservation of the total charge.  These two limits are distinct; the
hydrodynamic pole is the interpolation between them.

For a transverse vector source, \(k_iA_{ri}=0\) and \(A_{r0}=0\), the phase
decouples at quadratic order and
\[
  j^i=-\sigma\partial_tA_{ri}.
\]
With the electric-field convention
\eqref{eq:sk-hydro-constitutive-current-with-source}, this is Ohm response
\(j^i=\sigma E_i\).  The same \(\sigma\) therefore appears in the diffusion
constant, the vector-source response, and the noise coefficient below.

\section{Noise Positivity and Langevin Form}

The imaginary part of \eqref{eq:sk-hydro-diffusion-action} is
\begin{equation}
  \operatorname{Im} I_{\rm diff}
  =
  T\sigma\int dt\,d^dx\,B_{ai}B_{ai}.
  \label{eq:sk-hydro-imaginary-action}
\end{equation}
It is nonnegative precisely when \(\sigma\ge0\).  The same term has a
stochastic representation.  For each spatial component,
\begin{equation}
  \exp\!\left[
    -T\sigma\int B_{ai}B_{ai}
  \right]
  =
  \int D\xi\,
  \exp\!\left[
    -\int\frac{\xi_i\xi_i}{4T\sigma}
    +\ii\int B_{ai}\xi_i
  \right],
  \label{eq:sk-hydro-hubbard-stratonovich}
\end{equation}
with normalization independent of \(B_a\).  Thus the noise satisfies
\begin{equation}
  \langle \xi_i(t,\mathbf x)\xi_j(t',\mathbf x')\rangle
  =
  2T\sigma\,\delta_{ij}
  \delta(t-t')\delta^{(d)}(\mathbf x-\mathbf x').
  \label{eq:sk-hydro-noise-correlator}
\end{equation}
After the Hubbard--Stratonovich transformation, varying \(\varphi_a\) gives
the Langevin equation
\begin{equation}
  \partial_t n+\partial_i\bigl(-D\partial_i n+\sigma E_i+\xi_i\bigr)=0 .
  \label{eq:sk-hydro-langevin-diffusion}
\end{equation}
This derivation fixes the factor of \(2T\sigma\) in the noise correlator
from the same coefficient that governs the dissipative current.

\section{Dynamical KMS at Quadratic Order}

Assume the initial state is thermal with inverse temperature
\(\beta=1/T\) and that the long-wavelength effective action has the
classical dynamical KMS symmetry appropriate to the derivative order under
consideration.  In the diffusion sector, the transformation needed at this
order is
\begin{equation}
  B_{ai}(t,\mathbf x)
  \mapsto
  B_{ai}(-t,\mathbf x)
  +\ii\beta\,\partial_tB_{ri}(-t,\mathbf x),
  \label{eq:sk-hydro-classical-kms-ba}
\end{equation}
with the time-reversal parity of the spatial vector field included in
\(\partial_tB_{ri}(-t,\mathbf x)\).  To see the coefficient constraint, keep
one spatial component and write the dissipative-noise part as
\[
  L=-\sigma a x+\ii c a^2,
  \qquad
  a=B_{ai},
  \qquad
  x=\partial_tB_{ri}.
\]
Time reversal sends \(x\mapsto -x\), while the KMS shift sends
\(a\mapsto a+\ii\beta x\).  The transformed expression is
\[
  L'
  =
  \sigma ax+\ii\sigma\beta x^2
  +\ii c a^2-2c\beta ax-\ii c\beta^2x^2.
\]
Equality with \(L\), up to total derivatives and terms beyond the retained
order, requires the coefficient of \(ax\) to match:
\[
  \sigma-2c\beta=-\sigma.
\]
The \(x^2\) terms then cancel by the same condition.  Therefore
\begin{equation}
  c=\frac{\sigma}{\beta}=T\sigma,
  \label{eq:sk-hydro-kms-noise-coefficient}
\end{equation}
which is the coefficient in \eqref{eq:sk-hydro-diffusion-action}.  This is
the local effective-action form of the fluctuation--dissipation relation of
Chapter~\ref{ch:thermal-spectral-kubo-transport}.

\section{Relation to Stress-Tensor Hydrodynamics}

For stress-tensor hydrodynamics, the charge phase fields are replaced by
doubled maps from a fluid worldvolume into spacetime, together with doubled
metrics and, when charges are present, the phases above.  The structural
logic is the same:
\[
\begin{aligned}
  \text{gauge/diffeomorphism invariance}
  &\Longrightarrow
  \text{Ward identities},\\[2mm]
  I[A,A]=0
  &\Longrightarrow
  \text{closed-time-path normalization},\\[2mm]
  \operatorname{Im}I\ge0
  &\Longrightarrow
  \text{positive noise kernel},\\[2mm]
  \text{dynamical KMS}
  &\Longrightarrow
  \text{dissipation--fluctuation relations}.
\end{aligned}
\]
At first derivative order this reproduces the viscous stress tensor and
noise kernel used in Chapter~\ref{ch:thermal-hydrodynamic-fluctuations}.
At higher derivative order the classification becomes a local invariant
problem for doubled fields subject to the constraints above.  A term belongs
to the hydrodynamic effective action only after its frame dependence,
contact-term ambiguity, KMS transformation, and positivity properties have
been specified.

\section{The Microscopic Theorem Boundary}

The construction in this chapter is an effective-field-theory statement.  A
microscopic derivation from QFT would have to construct the real-time
generating functional in a thermal state, take the thermodynamic limit,
identify the hydrodynamic scaling limit, prove that the only nonanalytic
long-wavelength modes are the conserved densities, and show that the local
functional above is the asymptotic expansion of connected
Schwinger--Keldysh correlators in a controlled topology.  The local
positivity and KMS constraints are necessary consequences of the microscopic
closed time path; they are not by themselves a proof of thermalization.

\begin{openproblem}[Microscopic derivation of SK hydrodynamic action]
\label{op:microscopic-sk-hydrodynamic-action}
Starting from a continuum QFT with a KMS state, derive the hydrodynamic
Schwinger--Keldysh effective action as a controlled long-wavelength limit of
connected real-time correlators, including the topology in which the
derivative expansion is asymptotic, the treatment of contact terms, and the
relation between source-response kernels and commutator spectral functions.
\end{openproblem}
